% Part:sets-functions-relations
% Chapter: size-of-sets
% Section: enumerations-alt

\documentclass[../../../include/open-logic-section]{subfiles}

\begin{document}

\olfileid[pt]{sfr}{siz}{enm-alt}

\olsection{Enumerações e Conjuntos Enumeráveis}

\begin{editorial}
  Esta seção define enumerações como bijeções com (segmentos iniciais)
  de~$\Nat$, como se faz em teoria dos conjuntos. Por isso difere ligeiramente
  das definições em~\olref[enm]{sec} e repete os exemplos dessa seção. É também
  um pouco mais concisa.
\end{editorial}

Podemos especificar um conjunto finito simplesmente enumerando os seus
!!{element}s. Fazemo-lo quando definimos um conjunto assim:
\[
  A = \{a_1, a_2, \ldots, a_n\}.
\]
Supondo que os !!{element}s $a_1$, \dots, $a_n$ são todos distintos,
obtemos !!a{bijection} entre $A$ e os primeiros $n$ números naturais $0$, \dots,
$n-1$. Reciprocamente, como todo o conjunto finito tem apenas finitos
!!{element}s, todo o conjunto finito pode ser posto nessa correspondência. Em
outras palavras, se $A$ é finito, existe !!a{bijection} entre $A$ e $\{0, \dots,
n-1\}$, onde $n$ é o número de !!{element}s de~$A$.

Se admitirmos certos tipos de conjuntos infinitos, admitiremos também que alguns
conjuntos infinitos possam ser enumerados. Podemos precisar isto dizendo que um
conjunto infinito é enumerado por !!a{bijection} entre ele e todo o~$\Nat$.

\begin{defn}[Enumeração, no sentido conjuntista]
Uma \emph{enumeração} de um conjunto $A$ é !!a{bijection} cujo contradomínio é
$A$ e cujo domínio é ou um conjunto inicial de números naturais $\{0,
1, \ldots, n\}$ \emph{ou} todo o conjunto dos números naturais~$\Nat$.
\end{defn}

\begin{explain}
Há uma base intuitiva para este uso da palavra \emph{enumeração}. Dizer que
enumeramos um conjunto $A$ é dizer que existe !!a{bijection} $f$ que nos permite
\emph{contar} os elementos do conjunto $A$. O elemento de índice~$0$ é $f(0)$, o
de índice~$1$ é $f(1)$, \ldots o de índice~$n$ é $f(n)$\ldots.\footnote{Sim,
contamos a partir de~$0$. Claro que também poderíamos começar em~$1$. Não faria
grande diferença. Bastaria substituir~$\Nat$ por~$\PosInt$.} O raciocínio fica
ainda mais claro com o seguinte:
\end{explain}

\begin{defn}
  \ollabel{defn:enumerable}
  Um conjunto~$A$ é !!{enumerable} se e só se $A = \emptyset$ ou existe
  uma enumeração de~$A$. Dizemos que $A$ é !!{nonenumerable} se e só se $A$
  \emph{não} é !!{enumerable}.
\end{defn}

\begin{explain}
Logo um conjunto é !!{enumerable} se e só se é vazio ou se podemos usar uma
enumeração para contar os seus !!{element}s.
\end{explain}

\begin{ex}
Uma função que enumera os números naturais é simplesmente a função identidade
$\Id{\Nat} \colon \Nat \to \Nat$ dada por $\Id{\Nat}(n) = n$. Uma função que
enumera os números naturais \emph{positivos}, $\Nat^+ =
\Nat \setminus \{0\}$, é a função $g(n) = n + 1$, isto é, a função sucessora.
\end{ex}

\begin{prob}
Mostre que um conjunto $A$ é !!{enumerable} se e só se $A = \emptyset$ ou existe
!!a{surjection} $f\colon \Nat \to A$. Mostre que $A$ é !!{enumerable} se e só se
existe !!a{injection} $g\colon A \to \Nat$.
\end{prob}

\begin{ex}
As funções $f\colon \Nat \to \Nat$ e $g \colon \Nat \to \Nat$ dadas por
\begin{align*}
  f(n) & = 2n \text{ e}\\
  g(n) & = 2n+1
\end{align*}
enumeram, respectivamente, os números naturais pares e os ímpares. Mas nenhuma é
!!{surjective}, logo nenhuma é uma enumeração de~$\Nat$.
\end{ex}

\begin{prob}
Defina uma enumeração dos quadrados perfeitos $1$, $4$, $9$, $16$, \dots
\end{prob}

\begin{ex}
Seja $\lceil x \rceil$ a função \emph{teto}, que arredonda $x$ por excesso para
o inteiro mais próximo. Então a função $f \colon \Nat \to \Int$ dada por:
\[
  f(n) = (-1)^{n} \left\lceil\tfrac{n}{2}\right\rceil
\]
enumera o conjunto dos inteiros~$\Int$ do seguinte modo:
\[
\begin{array}{c c c c c c c c}
f(0) & f(1) & f(2) & f(3) & f(4) & f(5) & f(6) & \dots \\ \\
\big\lceil \tfrac{0}{2} \big\rceil & -\big\lceil \tfrac{1}{2}\big\rceil &  \big\lceil \tfrac{2}{2} \big\rceil & -\big\lceil \tfrac{3}{2} \big\rceil & \big\lceil \tfrac{4}{2} \big\rceil  & -\big\lceil \tfrac{5}{2}\big\rceil & \big\lceil \tfrac{6}{2} \big\rceil & \dots \\ \\
0 & -1 & 1 & -2 & 2 & -3 & 3& \dots
\end{array}
\]
Repare como $f$ gera os valores de $\Int$ ``saltitando'' entre inteiros
positivos e negativos. Também se pode ver $f$ definida por casos assim:
\[
f(n) = \begin{cases}
  \frac{n}{2} & \text{se $n$ é par}\\
  -\frac{n+1}{2} & \text{se $n$ é ímpar}
  \end{cases}
\]
\end{ex}

\begin{prob}
Mostre que se $A$ e $B$ são !!{enumerable}s, então $A \cup B$ também é.
\end{prob}

\begin{prob}
Mostre por indução em~$n$ que se $A_1$, $A_2$, \dots, $A_n$ são todos
!!{enumerable}s, então $A_1 \cup \dots \cup A_n$ também é.
\end{prob}

\end{document}
